\documentclass[11pt]{article}

\usepackage[margin=1in]{geometry}
\usepackage{amsmath, amssymb}
\usepackage{booktabs}
\usepackage{hyperref}
\usepackage{parskip}
\usepackage{titlesec}
\usepackage{xcolor}
\usepackage{enumitem}

\hypersetup{colorlinks=true}

\titleformat{\section}{\large\bfseries}{{\thesection.}}{0.5em}{}
\titleformat{\subsection}{\normalsize\bfseries}{{\thesubsection}}{0.5em}{}

\title{
  \textbf{STAT 405 Project Proposal} \\[0.4em]
  Hierarchical Bayesian Forecasting of U.S. State Unemployment Rates
}
\author{Calvin Du (Student ID: 16507378)}

\begin{document}
\maketitle

% Intro
\section{Team}

Calvin Du (Student ID: 16507378)

\section{Code Repository}

\url{https://github.com/calvingdu/stat405}

% questions 
\section{Scientific Questions}

As a STAT + ECON major, I am particularly interested in labour markets,
especially given the disruptions of the past several years including COVID-19,
the rise of AI, and recent political changes. These have all materially affected
employment patterns. Unemployment is probably the most consequential and
closely-watched labour market indicator signalling a good or bad economic environment.
State-level unemployment rates in the U.S.\ have huge variation mainly driven by 
local industry composition, geography, and sensitivity to national 
macroeconomic conditions (eg, California is more independent than other states). 
At the same time, federal monetary policy and typical economic growth cycles will always
affect all the states regardless of how independent they are. This creates a 
need to balance treating each state independently and pooling all states together.

\textbf{Primary question:}
\begin{quote}
Can I produce more accurate unemployment forecasts for each state by 
partially pooling information across the states, borrowing strength from 
the national trend rather than purely fitting each state by themselves?
\end{quote}

\textbf{Secondary question:}
\begin{quote}
Does including covariates like the federal monetary policy (like federal interest rates) 
and national GDP growth improve forecasts beyond an autoregressive model? 
\end{quote}

Both questions are naturally explored using Bayesian hierarchical models,
where the degree of pooling is itself inferred from the data.

% datasets
\section{Datasets}

\subsection{Primary Dataset: BLS State Unemployment}

The U.S.\ Bureau of Labor Statistics publishes monthly state-level
unemployment rates through its Local Area Unemployment Statistics (LAUS)
program. I've chosen to use this the "Employment status of the civilian noninstitutional population, seasonally adjusted" series
because it is the most widely cited and it's very clean and well-maintained. With it, I have 
a long time series of unemployment rates for all 50 states.

\textbf{Coverage:} I'll be using all 50 U.S.\ states, from January 2000 to December 2025,
monthly frequency (seasonally adjusted). This yields a $50 \times 312$
data matrix.

\textbf{URL:} \url{https://www.bls.gov/lau/rdscnp16.htm}

\medskip
\noindent\textbf{Dataframe structure (head):}

\begin{center}
\small
\begin{tabular}{llrrrrrr}
  \toprule
  \textbf{FIPS} & \textbf{State} & \textbf{Year} & \textbf{Month} 
    & \textbf{Labor force} & \textbf{Employed} & \textbf{Unemployed} & \textbf{Rate (\%)} \\
  \midrule
  01 & Alabama    & 1976 & 01 & 1,486,509 & 1,387,606 & 98,903  & 6.7 \\
  02 & Alaska     & 1976 & 01 &   160,709 &   149,226 & 11,483  & 7.1 \\
  04 & Arizona    & 1976 & 01 &   963,980 &   865,261 & 98,719  & 10.2 \\
  05 & Arkansas   & 1976 & 01 &   889,912 &   825,233 & 64,679  & 7.3 \\
  06 & California & 1976 & 01 & 9,770,553 & 8,873,133 & 897,420 & 9.2 \\
  \multicolumn{8}{c}{\ldots} \\
  \bottomrule
\end{tabular}
\end{center}

\medskip
\subsection{In order to answer the secondary question, I will also pull in national macroeconomic covariates from the Federal Reserve Economic Data}
For background, FRED is a widely used database of economic time series maintained by the Federal Reserve Bank of St.\ Louis. 
It's extremely useful for economic research and forecasting, and is used a lot by economists and data scientists. 

The data I'll be using are:
\begin{itemize}[noitemsep]
  \item \textbf{Federal funds effective rate (\texttt{FEDFUNDS}) - Monthly}: This is the interest rate at which depository institutions trade federal funds with each other. It is a key tool of U.S.\ monetary policy that has a strong influence on economic activity and unemployment.
  \item \textbf{Real Gross Domestic Product (\texttt{A191RL1Q225SBEA}) - Quarterly}: This is the real gross domestic product, which is a measure of overall economic output. The growth rate of real GDP is closely related to unemployment.
\end{itemize}

\medskip
\noindent Both the series are downloaded as plain CSV files directly from FRED 
(\url{https://fred.stlouisfed.org}) and cover January 2000 to December 2025.

\medskip

\begin{center}
\small
\begin{tabular}{lr}
  \toprule
  \textbf{observation\_date} & \textbf{FEDFUNDS} \\
  \midrule
  2000-01-01 & 5.45 \\
  2000-02-01 & 5.73 \\
  \multicolumn{2}{c}{\ldots} \\
  2025-12-01 & 3.72 \\
  \bottomrule
\end{tabular}
\hspace{1.5em}
\begin{tabular}{lr}
  \toprule
  \textbf{observation\_date} & \textbf{A191RL1Q225SBEA} \\
  \midrule
  2000-01-01 & 1.5 \\
  2000-04-01 & 7.5 \\
  \multicolumn{2}{c}{\ldots} \\
  2025-10-01 & 0.7 \\
  \bottomrule
\end{tabular}
\end{center}

\subsection{Backup Dataset: OECD Country Unemployment}

As a backup, the Organisation for Economic Co-operation and Development (OECD) publishes
unemployment data across various countries. An alternative to only focusing on U.S states is to
apply the same hierarchical modelling approach to a broader set of regions across multiple countries.
I could look at unemployment rates in an Asian country and bring in macroeconomic covariates from that region 
(such as Japan unemployment rate taking into account China and South Korea rates, or how allies of Japan doing). 

The URL for this data is: \\
\url{https://sdmx.oecd.org/public/rest/data/OECD.SDD.TPS,DSD_LFS@DF_IALFS_UNE_M,1.0/AUS+AUT+BEL+CAN+CZE+DNK+FIN+FRA+DEU+GRC+HUN+IRL+ITA+JPN+KOR+NLD+NZL+NOR+POL+PRT+SVK+SVN+ESP+SWE+CHE+TUR+GBR+USA..PT_LF_SUB._Z.Y._T.Y_GE15..M?startPeriod=2000-01&dimensionAtObservation=AllDimensions&format=csvfile
}

% methods
\section{Methodology}

\subsection{Model 1: Independent AR(1) (Naive)}

Each state is modelled independently with no pooling. For state $s$ and
time $t$:

\begin{align*}
  \alpha_s &\sim \mathcal{N}(5,\; 3) & (\text{baseline level of unemployment is around 5\%}) \\
  \beta_s  &\sim \mathcal{N}(0.5,\; 0.3) & (\text{baseline autoregressive coefficient - moderate persistence}) \\
  \sigma   &\sim \mathrm{Exp}(1)  & (\text{shared SD across states}) \\[4pt]
  y_{s,t} \mid \alpha_s, \beta_s, \sigma
           &\sim \mathcal{N}\!\left(\alpha_s + \beta_s\, y_{s,t-1},\; \sigma\right)
\end{align*}

This serves as the naive baseline. Since states are fit independently,
estimates for small or volatile states will be noisy, with no
regularisation from the national trend or neighbouring states. 
This model does allow states whose unemployment rates genuinely
differ from the national picture to have their 
own parameters without being pulled toward the national mean.

\subsection{Model 2: Hierarchical AR(1) with Covariates}

State parameters are drawn from national priors, enabling
partial pooling. Macroeconomic covariates enter the likelihood as shared
predictors:

\begin{align*}
  \mu_\alpha &\sim \mathcal{N}(5,\; 3) 
    & \text{(national mean intercept)} \\
  \sigma_\alpha &\sim \mathrm{Exp}(1/2) 
    & \text{(spread of intercepts across states)} \\
  \mu_\beta &\sim \mathcal{N}(0.5,\; 0.3) 
    & \text{(national mean persistence coefficient)} \\
  \sigma_\beta &\sim \mathrm{Exp}(1/0.3) 
    & \text{(spread of persistence coefficients across states)} \\
  \boldsymbol{\gamma} &\sim \mathcal{N}(0,\; 1) 
    & \text{(national economic covariate coefficients)} \\
  \sigma &\sim \mathrm{Exp}(1) 
    & \text{(shared SD across states)} \\[6pt]
  \alpha_s \mid \mu_\alpha, \sigma_\alpha &\sim \mathcal{N}(\mu_\alpha,\; \sigma_\alpha)
    & \text{(state intercept, partially pooled)} \\
  \beta_s \mid \mu_\beta, \sigma_\beta &\sim \mathcal{N}(\mu_\beta,\; \sigma_\beta)
    & \text{(state persistence coefficient, partially pooled)} \\[6pt]
  y_{s,t} \mid \alpha_s, \beta_s, \boldsymbol{\gamma}, \sigma
    &\sim \mathcal{N}\!\left(
        \alpha_s + \beta_s\, y_{s,t-1} + \boldsymbol{\gamma}^\top \mathbf{x}_t,\;
        \sigma
      \right) 
\end{align*}

where $\mathbf{x}_t = (\texttt{FEDFUNDS}_t,\; \texttt{GDP\_growth}_t)^\top$
contains two national covariates, with
\texttt{GDP\_growth} interpolated from quarterly to monthly.
Unlike Model 1, $\alpha_s$ and $\beta_s$ are tied together through shared
hyperparameters $(\mu_\alpha, \sigma_\alpha, \mu_\beta, \sigma_\beta)$ which represent
the national distribution of state-level parameters. This allows for
shrinkage: states with less data or more noise will be pulled toward the national mean, while states with strong signals can deviate more.
The key variable is $\sigma_\beta$ where a small value forces all states toward the same estimated national level
coefficient $\mu_\beta$, while a large value allows states to behave
independently. 

\subsection{Posterior Inference}

\begin{itemize}
  \item \textbf{Method 1: SMC via simPPLe:} \\
        SMC will be a good fit for the independent model where parameters are uncorrelated across states,
        and it will be interesting to see how it performs on the hierarchical model where parameters are correlated and more complex.
        However, it may struggle with the large number of parameters and the hierarchical structure, leading to slower convergence and less accurate estimates. 
  \item \textbf{Method 2: MCMC/HMC via Stan} \\
        MCMC is more flexible for complex hierarchical models and can handle the correlated parameters better. 
        Also it will be faster to handle so many parameters and more efficient for the hierarchical model.
\end{itemize}

The comparison will examine:
\begin{itemize}
  \item Posterior mean and credible interval and agreement between methods
  \item Computational cost including runtime and effective sample size (ESS) per second
  \item MCMC diagnostics such as $\hat{R}$, ESS
\end{itemize} 


% comparison
\subsection{Model Comparison and Evaluation}

The data are split like any timeseries: January 2000 to December 2023 will be set as training data, and
January 2024 to December 2025 (24 months) will be the test set.
Evaluation metrics:
\begin{itemize}[noitemsep]
  \item  RMSE of the posterior predictive mean per state
  \item Shrinkage plot: posterior mean $\hat{\beta}_s$ (Model 2) vs.\ MLE
        from independent fits (Model 1)
  \item Prior sensitivity: re-fit Model 2 with wider/narrower priors on
        $\sigma_\beta$
\end{itemize}

\subsection{Diagnostics and Validation}

\begin{itemize}[noitemsep]
  \item $\hat{R}$ for all parameters
  \item Trace plots for hyperparameters $(\mu_\beta, \sigma_\beta)$ and
        selected state-level parameters
  \item Non-centred parameterisation if divergences arise in Stan
\end{itemize}

% misspecifcations
\section{Potential Model Misspecifications}

\begin{itemize}
  \item \textbf{COVID-19:} unemployment spiked to $\sim$15\%
        nationally during 2020 and is an extreme outlier. A sensitivity analysis excluding 2020 will be reported.
  \item \textbf{Structural breaks:} the relationship between unemployment and covariates may change over time, especially with recent economic disruptions. 
  A time-varying coefficient model could be explored as an extension.
\end{itemize}


\end{document}
